Learning-guided metaheuristics are routinely reported to beat their non-learning counterparts, yet the comparisons seldom separate what the learned policy adds from what the search structure, the CPU budget and the tuning already provide. We propose a budget-fair attribution protocol (a random control at a matched budget; equal CPU time, with the gain split into a decision effect and a cost effect; a static competitor validated offline; a cross-fitted oracle headroom whose bias we characterise; power checks; hypotheses recorded before each run) and apply it to Iterated Greedy (IG) for parallel machine scheduling, with a permutation flowshop and a genetic algorithm as further hosts, in \nbatteries{} registered batteries (\npredheld{} of \npredtotal{} point predictions held). Four findings. (1) The destruction--reconstruction mechanism, not the policy, explains the IG's advantage over an iterated local search; the share of a local search's gain that is only budget ranges from $28\%$ to $93\%$ with the implementation. (2) At equal iterations no online learner beats random choice by more than $0.04$ points on the synthetic classes. A validated static action, or a mixture where actions are complementary, was never beaten beyond the equivalence margin by an online learner on any host, including three published families of unrelated-machine instances and the standard benchmark of Fanjul-Peyro and Ruiz ($700$ test instances, $147$ learner--cell pairs, and again with a ten times longer horizon), nor by an irace-configured algorithm, two published adaptive-operator-selection designs or a neural policy trained offline. (3) The oracle headroom net of run-to-run noise is a necessary condition for Q-Learning to pay, not a sufficient one. (4) The apparent exceptions come from the learner's implementation: the standard Q-Learning repeats the action listed first (zero-initialised table, ties broken by index), so its gains and losses, and its apparent advantage at equal CPU time, measure where a good action stood in the list; an order-invariant version gains beyond the margin in none of $69$ cells of four collections of instances (on the held-out published instances of $P||\Cmax$ it loses where the standard version gained), and the protocol therefore requires reporting a learner's order sensitivity. What headroom does reward is knowledge acquired before the run: a neural policy trained offline recovers about two thirds of the validated gain on the standard benchmark, $3.7$ times what the best online learner finds, although most of that is a constant mixture of actions.
